\BeleskeID{01-s009}

\section*{Objašnjenje iz originala}

Izvor: predavanje 01, PDF strana 5, gornji slajd.

% element: 01-s009:e01

Rednim i paralelnim vezivanjem prekidača mogu se ostvariti kombinacije logičkih uslova koje opisuju rad realnog sistema.

% element: 01-s009:e02

Lift polazi sa trećeg sprata; prima poziv sa drugog sprata za kretanje nadole i sa četvrtog sprata za kretanje nagore. Primer zatim zahteva pokretanje nagore, zaustavljanje na četvrtom spratu i otvaranje vrata. Ovo je opis problema, a ne kompletan algoritam upravljanja.

% element: 01-s009:e03

Svakoj tvrdnji pridružuje se promenljiva. A označava tvrdnju da je lift na trećem spratu, B poziv sa drugog sprata za kretanje nadole, a Y komandu za pokretanje nagore. Elipsa u funkciji označava dodatne uslove. Funkcija se može realizovati odgovarajućom šemom prekidača.

\section*{Dopunsko objašnjenje}

% element: 01-s009:d01

Nije moguće odrediti konkretnu funkciju Y samo iz tri navedene promenljive; ostali ulazi i pravila odlučivanja nisu dati u primeru.
